% =================================================================
% Part XV: Universal Adapter Framework
%   NRMO/Loom を domain 非依存にする adapter 構造.
%   Hybrid (高出力 × 破綻回避) + 4 domain 実証 + C++/Python/Frontend.
% =================================================================

\part{Universal Adapter Framework}
\label{part:universal_adapter}

\chapter{Domain Independence of NRMO}
\label{ch:domain_independence}

\section{From Civilisation to Any Domain}

Parts I--XIV developed NRMO/Loom primarily against abstract benchmark
worlds and the civilisation simulation. This part establishes a
stronger claim: the core is \emph{domain-independent}. The same
behavioral core (Loom v3.1), the same Safety Floor, the same Sociable
Shadow, and the same Hybrid composition apply to any sequential
decision problem that contains an absorbing failure state.

The defining property of NRMO's applicability is not ``civilisation''
but the structure
\begin{equation}
\max_{\pi}\ \mathbb{E}[U(\pi)]
\quad\text{s.t.}\quad
\Pr[s_t \in \mathcal{R}\ \text{for some } t] \le \varepsilon,
\end{equation}
where $\mathcal{R}$ is a set of absorbing ruin states. Financial
portfolios (margin call), retail operations (insolvency), health
trajectories (chronic illness), and careers (irreversible dead-ends)
all share this structure.

\section{The 6D State as a Cross-Domain Language}

The single mechanism that makes this work is the shared
six-dimensional state $S=(R,E,G,O,K,X)$. Every domain answers the same
six questions:

\begin{center}
\begin{tabular}{cll}
\toprule
Axis & Generic meaning & Example (finance) \\
\midrule
$R$ & resource / liquidity (safe$\uparrow$) & portfolio value, cash \\
$E$ & sustainability        & return stability (low vol) \\
$G$ & governance / discipline & diversification \\
$O$ & optionality           & cash ratio / rebalance room \\
$K$ & knowledge             & trend confidence \\
$X$ & exposure (danger$\uparrow$) & drawdown / volatility \\
\bottomrule
\end{tabular}
\end{center}

\chapter{The DomainAdapter Interface}
\label{ch:domainadapter}

\section{Required and Optional Methods}

A new domain is connected by implementing a thin adapter. Four methods
are required; one is optional (for Hybrid mode).

\begin{description}[leftmargin=1.6em]
  \item[\texttt{to\_loom\_state(ds)}] map domain state to $(R,E,G,O,K,X)$.
  \item[\texttt{apply\_action(action, ds, rng)}] apply a Loom action and
    return the next domain state.
  \item[\texttt{is\_ruin(ds)}] test whether the state is an absorbing
    failure.
  \item[\texttt{compute\_reward(prev, next)}] scalar feedback.
  \item[\texttt{propose\_high\_output(ds, rng)} \rm(optional)] a
    high-output action proposal used by the Hybrid controller.
\end{description}

The controller, the Safety Floor, the sparse-activation logic, and the
Sociable Shadow are all shared and never re-implemented per domain.

\section{The Hybrid Controller}

The composition that unifies high output and non-ruin is:

\begin{center}
\begin{tikzpicture}[
  node distance=7mm,
  box/.style={draw,rounded corners,minimum width=4.2cm,minimum height=8mm,font=\small,align=center},
]
\node[box,fill=blue!8]  (hi)   {High-output proposal (domain)};
\node[box,fill=green!8,below=of hi]  (ctrl) {Loom Control (detect + risk)};
\node[box,fill=red!8,below=of ctrl]  (sf)   {Safety Floor (recover-first)};
\node[box,fill=yellow!12,below=of sf](act)  {domain action};
\node[box,fill=gray!10,right=14mm of ctrl] (sh) {Sociable Shadow\\(trace only)};
\draw[-{Latex}] (hi)--(ctrl);
\draw[-{Latex}] (ctrl)--(sf);
\draw[-{Latex}] (sf)--(act);
\draw[-{Latex},dashed] (ctrl)--(sh);
\end{tikzpicture}
\end{center}

In calm states the high-output proposal passes through unchanged; in
dangerous states the Safety Floor substitutes a \emph{recover-first}
action --- importantly, boosting governance/repair rather than merely
cutting growth, since naive growth-cutting itself triggers a
resource-floor collapse (an error observed and corrected during
development).

\chapter{Empirical Results across Four Domains}
\label{ch:four_domains}

Each domain uses the identical core; only the adapter differs. The
``Hybrid effect'' column compares the Hybrid controller against the
conservative Loom-only baseline.

\begin{center}
\begin{tabular}{lll}
\toprule
Domain & Ruin condition & Hybrid effect (vs Loom-only) \\
\midrule
Civilisation   & resource / exposure collapse & cum.\ prod $\approx$98--100\% of high-output, ruin 0\% \\
Finance        & drawdown $>40\%$             & max drawdown $18.6\%\!\to\!1.1\%$ \\
Retail (store) & cash $<15$                   & survival $35\!\to\!85$--$108$ steps \\
Health         & disease risk $>80$           & score $973\!\to\!1134$, ruin 0\% \\
\bottomrule
\end{tabular}
\end{center}

\paragraph{Cross-language validation.} The C++ store demo independently
reproduces the trend (Hybrid revenue $1221.5$ vs Loom-only $231.0$),
confirming the effect is a property of the control structure rather
than of a particular language or RNG.

\paragraph{Regularities.}
(i) Loom-only is conservative in every domain. (ii) The Hybrid
consistently dominates it. (iii) The adapter's dynamics model is a
separate concern from the core --- a harsh domain model can still
produce ruin, but that reflects domain modelling, not a controller
failure.

\chapter{Reference Implementations}
\label{ch:reference_impl}

\section{Python (core interface)}
\begin{lstlisting}[language=Python]
class DomainAdapter(ABC):
    name: str
    state_semantics: dict   # meaning of R,E,G,O,K,X
    ruin_condition: str
    @abstractmethod
    def to_loom_state(self, ds): ...      # -> WorldState(R,E,G,O,K,X)
    @abstractmethod
    def apply_action(self, action, ds, rng): ...
    @abstractmethod
    def is_ruin(self, ds) -> bool: ...
    @abstractmethod
    def compute_reward(self, prev, nxt) -> float: ...
    def propose_high_output(self, ds, rng):   # Hybrid (optional)
        return Action("invest", "C")

# connect a new domain in three lines:
register_adapter("store", StoreOperationAdapter)
ctrl = make_controller("store", seed=42, use_hybrid=True)
result = ctrl.run_episode(StoreState(), horizon=150, rng=rng)
\end{lstlisting}

\section{C++ (header-only core)}
\begin{lstlisting}[language=C++]
template <typename DomainState>
Action LoomController::decide(const DomainAdapter<DomainState>& adapter,
                              const DomainState& ds) {
    WorldState s = adapter.to_loom_state(ds);
    double risk  = adapter.risk_proximity(ds);
    LoomMode mode = select_mode(s, risk);
    Action a = (use_hybrid_ && risk < EMERGENCY_RISK)
             ? adapter.propose_high_output(ds, rng_)   // high output
             : conservative_action(mode);              // Loom safe
    return apply_safety_floor(a, s);                   // final gate
}
\end{lstlisting}

\section{Browser Front-End}
A single-file HTML console ports the core decision logic to JavaScript.
With no programming, a user selects a domain, chooses Hybrid or
Loom-only, sets horizon and seed, runs an episode, and inspects the
6D state trajectory, the score trajectory, and the mode-usage
histogram. It runs entirely offline.

\chapter{Significance for NRMO}
\label{ch:adapter_significance}

The Universal Adapter Framework completes a conceptual arc. NRMO began
as a theory of survival-first decision making; the Loom cycle produced
a robust behavioral core; the Sociable work clarified that detection
and intervention must be separated; the Hybrid showed that safety and
output are composable rather than opposed. This part shows that the
result is not tied to any one domain:

\begin{equation*}
\boxed{\;
\text{NRMO}_{\text{current}}
= \text{Hybrid}\big(\text{high-output engine},\ \text{Loom v3.1}\big)
+ \text{Sociable Shadow}
\;+\;\text{DomainAdapter}
\;}
\end{equation*}

Any problem with an absorbing failure state can join by answering six
questions (the $R,E,G,O,K,X$ semantics) and four method signatures.
``High output under a non-ruin guarantee'' is thereby made portable.
